\subsection{Algorithme de déplacement}

% -----------------------------PSEUDO CODE----------------------------------------------
\noindent\textbf{procedure} ROBOT\_evolutionRobot ()

\begin{tabular}{ll}
    \textbf{Déclaration:} & prochaineAction: \textbf{Chaine de caractères}\\
     & etat: \textbf{ROBOT\_EtatDAvancement}\\
\end{tabular}

\noindent\textbf{début}

    \begin{algorithm}[H]
        etat $\leftarrow$ Avancer\\
        \Tq{prochaineAction $\neq$ FIN}{
            \textbf{lire} (prochaineAction)\\
            \Suivant{etat}{
                \Cas{Avancer}{ROBOT\_avancer (etat)}
                \Cas{Intersection}{ROBOT\_intersection (etat, prochaineAction)}
                \Cas{Urgence}{ROBOT\_urgence ()}
                }
        }
    \end{algorithm}

    \noindent\textbf{fin}
% -----------------------------PSEUDO CODE----------------------------------------------
\begin{center}
    \textit{Algorithme de ROBOT\_evolutionRobot}
\end{center}

Pour élaborer l'algorithme de déplacement du robot nous avons considéré qu'il pouvait se retrouver dans 4 états différents:
\begin{itemize}
    \item En avancée
    \item Sur une intersection
    \item En situation d'urgence
    \item Sorti du labyrinthe
\end{itemize} 

Nous avons donc une variable $etat$ du type \textbf{ROBOT\_EtatDAvancement} qui est mise à jour à chaque tour d'une boucle qui ne s'arrête que lorsque la suite d'ordre que nous lui fournissons s'achève.
Selon la valeur de cette variable $etat$, différentes fonctions sont appelées afin de faire évoluer le Robot.

La première est la fonction ROBOT\_avancer, qui met les moteurs en marche avant et qui adapte la trajectoire du robot grâce aux capteurs.

Ensuite, la fonction ROBOT\_intersection agit selon l'ordre lu par l'algorithme principal. Elle fait tourner le Robot à droite ou à gauche ou bien le fait continuer tout droit.
Le Robot effectue un virage sur place, nous avons donc paramétré un temps d'attente afin de positionner correctement le centre de rotation du robot sur l'intersection.

La fonction, ROBOT\_urgence, est la fonction qui est appelée lorsqu'un obstacle est détecté. Son rôle est d'arrêter les moteurs et d'allumer le buzzer.
